% Nouveau EGA V (cote 25): archive cote 25, site page 183 (PDF page 184).
% typescript. Draft compared: none.
% AI-transcribed by Claude Opus 5.5 (claude-opus-5-5, adaptive thinking, effort high), Message Batches API, batch ega25b (msgbatch_01A7n5YuXubUPNK1zsepi8RG), 2026-11-13; not yet checked by eye.
\scanpage{25/183}
\mspage{\hand{10}.}

ments de base. Des arguments purements formels,\xout{\illegible{2}} joints à la
``théorie de la descente plate'' de \inserted{EGA} V \unsure{1|I} , montreront alors que le foncteur
$(\underline{E},Q) \hand{$\rightsquigarrow$} P$ établit une \uline{équivalence} de la c atégorie fibrée des formes
quadratiques lisses sur des Modules localement libres de rang n, de base
un préschéma S quelconque, avec la catégorie fibrée des \uline{fibrés principaux}
\uline{homogènes} sur un préschéma quelconque S, \uline{de groupe} $\underline{O}(n)_S$ .

16.13. Soit Q une forme quadratique lisse sur un $\underline{E}$ loc libre de
rang fini, alors on vérifie facilement que pour $s \hand{$\in$} S$, $\underline{O}(Q)$ est lisse
sur un voisinage ouvert de s , \xout{\illegible{1} k(s) \illegible{1} de \illegible{1}} sauf exactemnt dans
le cas où k(s) est de car 2 et le rang de $\underline{E}$ en s est impair. Supposons
que cette circonstance ne se produi\xout{\unsure{se}} pour aucun $s \hand{$\in$} S$, alors (et alors
seulement) le P envisagé plus haut est lisse sur S (étant un fibré
principal homogène sous $\underline{O}(n)_S$ , donc lisse sss ce dernier l'est). Utili\xout{\illegible{1}}
sant le ``lemme de Hensel'' , il s'ensuit que \xout{\illegible{1}} l'on peut alors aux
conditions équivalentes de 16.9. adjoindre la condition \xout{\unsure{ex}} équivalente

(iii ter) \xout{\illegible{3}} Comme (iii), en supposant $S' \hand{$\to$} S$
\uline{étale} surjectif.

En fait, il y a mieux (en gardant toujours les hypothèses précéden-%
tes), car il résulte de la théorie générale des schémas en groupes réductifs
et fibrés principaux homogènes sur iceux (cf SGAD XXIV) que \xout{\illegible{1}} si Q est
lisse, alors tout point s admet un voisinage ouvert U, et un morphisme
étale surjectif \uline{fini} $S' \hand{$\to$} U$, tel que $Q_{S'}$ soit la forme standart. Si par exem
% NOTE: right edge of the page cut; the end of "exem-" is not visible
ple S est local, on peut dans (iii ter) supposer en plus $S' \hand{$\to$} S$ \uline{fini}.

\inserted{16.14.} \hand{(}\over{L}{C}es résultats \inserted{signalés dans 16.13.} \hand{(}à commencer par la condition (iii ter) pour une Q
lisse)\hand{)} tombent en défaut quand on abandonne l'hypothèse supplémentaire
de rang et de car. , par exemple lorsque S=Spec(k), k corps \inserted{imparfait} de car.2 , et
que E est de rang 1, car la forme quadratique at$^{2}$ , pour $\hand{$a \in$} k - k^{2}$ , ne peu
% NOTE: right edge of the page cut after "peu"

\margin{183}
